% Werke X/1 (1917), native191–195, printed184–188.
\begingroup
\setlength{\headheight}{25pt}
\makeatletter\tagsleft@true\let\veqno\@@leqno\makeatother
\SourceLeaf{191}\NoSourceHead
\SourceGroup{IV}
\HistoricTitle{[INVESTIGATIO FUNCTIONUM QUAE EX EVOLUTIONE MEDIORUM\\ ARITHMETICO-GEOMETRICORUM ORIUNTUR.]}
\HistoricCitation{[Auszüge aus der Scheda Ac. Varia, Novbr. 1799.\\
Imprimis de Integrali \(\displaystyle\int\frac{\partial\varphi}{\surd(1+\mu\mu\sin\varphi^2)}\). 22. XI. S. 7 ff.]}
\NumberedItem{1}
\noindent [S. 7]
\begin{equation}\tag*{[1]}\frac{\varpi}{\pi}
=\left(1+\left(\frac12\right)^2\frac12
+\left(\frac12.\frac34\right)^2\frac14
+\left(\frac12.\frac34.\frac56\right)^2\frac18\ldots\right)\surd\frac12\end{equation}
{\small
\begin{equation}\tag*{[2]}\begin{aligned}
=\biggl(1&+\frac{1.2.3}{2.4.4}.\frac19
+\frac{1.2.3.5.6.7}{2.4.4.6.8.8}.\frac1{81}\\
&+\frac{1.2.3.5.6.7.9.10.11}{2.4.4.6.8.8.10.12.12}.\frac1{729}\ldots\biggr)\surd\frac23
\end{aligned}\end{equation}
\begin{equation}\tag*{[3]}\begin{aligned}
=\biggl(1&+\frac14.\frac34.\frac19+\frac14.\frac34.\frac58.\frac78.\frac1{81}\\
&+\frac{1.3.5.7.9.11}{4.4.8.8.12.12}.\frac1{729}\ldots\biggr)\surd\frac23.
\end{aligned}\end{equation}}
\noindent [S. 9]
{\small
\begin{equation}\tag*{[4]}\begin{aligned}
\frac{\varpi}{\pi}-\frac2{\varpi}
=\biggl(&\frac12.\frac13+\frac12.\frac34.\frac34.\frac56.\frac1{27}\\
&+\frac12.\frac34.\frac34.\frac56.\frac56.\frac78.\frac9{10}.\frac1{243}+\text{etc.}\biggr)\surd\frac16
\end{aligned}\end{equation}
\begin{equation}\tag*{[5]}\begin{aligned}
\frac{\varpi}{\pi}-\frac2{\varpi}
=\biggl(&\frac14.\frac13+\frac14.\frac34.\frac58.\frac1{27}
+\frac14.\frac34.\frac58.\frac78.\frac9{12}.\frac1{243}\\
&+\frac14.\frac34.\frac58.\frac78.\frac9{12}.\frac{11}{12}.\frac{13}{16}.\frac1{2187}\ .\,.\biggr)\sqrt{\frac23}.
\end{aligned}\end{equation}}
\ClosingRule
\clearpage
\SourceLeaf{192}\ArithmeticHead{185}{R}{\scriptsize EVOLUTIO MEDIORUM ARITHMETICO-GEOMETRICORUM.}
\NumberedItem{2}
\noindent [S. 9]
\begin{equation}\tag*{[6]}\frac1{\surd(1-cz)}=A+A'c+A''c2+A'''c3+\text{etc.}\end{equation}
{\small
\[\begin{aligned}
A&=1+\frac{1.3}{4.4}zz+\frac{1.3.5.7}{4.4.8.8}z^4+\frac{1.3.5.7.9.11}{4.4.8.8.12.12}z^6+\text{etc.}\\
A'&=2z\left(\frac14+\frac{1.3.5}{4.4.8}zz+\frac{1.3.5.7.9}{4.4.8.8.12}z^4+\frac{1.3.5.7.9.11.13}{4.4.8.8.12.12.16}z^6+\text{etc.}\right)\\
A''&=2z^2\left(\frac{1.3}{4.8}+\frac{1.3.5.7}{4.4.8.12}zz+\frac{1.3.5.7.9.11}{4.4.8.8.12.16}z^4+\text{etc.}\right)\\
A'''&=2z^3\left(\frac{1.3.5}{4.8.12}+\frac{1.3.5.7.9}{4.4.8.12.16}zz+\frac{1.3.5.7.9.11.13}{4.4.8.8.12.16\ 20}z^4+\text{etc.}\right)\\
&\text{etc.}
\end{aligned}\]}
\[\begin{aligned}
2zA-4A'+3zA&=0\\
3zA'-8A''+5zA'''&=0\\
5zA''-12A'''+7zA^{\mathrm{IV}}&=0\\
\text{etc.}
\end{aligned}\]
\[\frac{A'}A=\cfrac{2z}{4-\cfrac{9zz}{8-\cfrac{25zz}{12-\cfrac{49zz}{16-\text{etc.}}}}}\]
\noindent Terminus constans duarum expressionum
\[\begin{gathered}
(4+4z+zz\cos\varphi^2)^{-\frac12}\\
(4+(8z+4zz)\cos\varphi^2)^{-\frac12}
\end{gathered}\]
\noindent idem est et quidem
\[\frac1{2\ \text{Term. med. Ar. G. inter 1 et }1+z}.\]
\par\vfill\noindent X1.\hfill24
\clearpage
\SourceLeaf{193}\ArithmeticHead{186}{L}{ANALYSIS.\quad NACHLASS.}
\NumberedItem{3}
\noindent [S. 10]
\par\indent Medium arithmetico-geometricum inter \(1\) et \(x\) per huius modi formulam exhiberi potest, apprime utilem quoties \(x\) est numerus satis magnus
\begin{equation}\tag*{[7]}Mx=\frac{C(x-\alpha x^{-1}-\beta x^{-3}-\gamma x^{-5}-\text{etc.})}
{\log.(4x-ax^{-1}-bx^{-3}-cx^{-5}-\text{etc.})}.\end{equation}
\noindent Hic factor constans
\[C=\frac12\pi=1{,}57079632\ldots;\]
\[\alpha=\frac14;\ \beta=\frac5{64};\ \gamma=\frac{11}{256};\
\delta=\frac{469}{36384};\ \varepsilon=\frac{1379}{2^{16}}\ldots;\]
\[a=1;\ b=\frac9{32};\ c=\frac{19}{128};\ d=\frac{4769}{49152};\]
{\small\begin{equation}\tag*{[8]}\text{Denominator}=\log.4x-\frac14x^{-2}-\frac{13}{128}x^{-4}-\frac{23}{384}x^{-6}-\frac{2701}{2^{16}}x^{-8}\end{equation}}
\ClosingRule
\begin{equation}\tag*{[9]}\begin{gathered}
\text{Obs[ervatio]: Numerator, omisso factore }C,\text{ pro }x=2,\text{ fit}\\
=M3=1{,}863617.
\end{gathered}\end{equation}
\par\indent Posito denominatore \(=\log.z\) erit
{\small\begin{equation}\tag*{[10]}\left\{\begin{aligned}
x={}&\frac14z+z^{-1}+\frac12z^{-3}-2z^{-5}-\frac14z^{-7}-3z^{-9}-\frac12z^{-11}+6z^{-13}\\
&+\frac34z^{-15}-22z^{-17}-3\frac12z^{-19}+64z^{-21}\\
={}&\frac1{\frac4z-\frac{16}{z^3}+\frac{56}{z^5}-\frac{160}{z^7}}.
\end{aligned}\right.\end{equation}}
\par\indent Posito numeratore \(=Cy\) erit
\begin{equation}\tag*{[11]}x=y+\frac14y^{-1}+\frac1{64}y^{-3}-\frac1{256}y^{-5}-\frac{23}{16384}y^{-7}\end{equation}
\noindent Denominator fit pro \(x=1+t\)
\clearpage
\SourceLeaf{194}\ArithmeticHead{187}{R}{\scriptsize EVOLUTIO MEDIORUM ARITHMETICO-GEOMETRICORUM.}
\begin{equation}\tag*{[12]}=-\frac{A}{\log.\frac18t-\frac1{16}tt+\frac5{128}t^3-\frac7{256}t^4+\frac{337}{16384}t^5},\end{equation}
\[A\text{ proxime}=4{,}9348=\frac12\pi\pi.\]
\NumberedItem{4}
\noindent [S. 11]
\begin{equation}\tag*{[13]}\text{Posito numeratore}=\frac CQ\end{equation}
\noindent erit
\begin{equation}\tag*{[14]}Q=\frac1x+\left(\frac12\right)^2x^{-3}
+\left(\frac{1.3}{2.4}\right)^2x^{-5}
+\left(\frac{1.3.5}{2.4.6}\right)^2x^{-7}+\text{etc.}\end{equation}
\begin{equation}\tag*{[15]}Q=\text{Pars ipsius }(xx-\cos\varphi^2)^{-\frac12},\text{ quae a }\varphi\text{ est libera}\end{equation}
\begin{equation}\tag*{[16]}\begin{gathered}
=\frac1{\frac12\pi}\times\text{valor integralis}\\
\int\surd\frac{1-rr}{xx-rr}dr\qquad\int\frac{dr}{\surd(1-rr)(xx-rr)}\\
\text{ab }r=0\text{ usque ad }r=1\\
\text{spectando }x\text{ tamquam constans}
\end{gathered}\end{equation}
\begin{equation}\tag*{[17]}\text{Ponendo }\frac12\frac{\surd x-1}{\surd x+1}=q\end{equation}
\noindent erit numerator
\begin{equation}\tag*{[18]}=-A:\log.(q+2q^5+15q^9+150q^{13}\end{equation}
\begin{equation}\tag*{[19]}=-A:\log.\varphi q.\end{equation}
\noindent Coefficientes seriei deducuntur ex aequatione
\[\left(\varphi\surd\frac{q}{1+4qq}\right)^2=\varphi q\]
{\small\[\left(\text{Forsan }\varphi q=q+\frac{1.2}{1.1}q^5+\frac{1.2.5.6}{1.1.2.2}q^9+\frac{1.2.5.6.9.10}{1.1.2.2.3.3}q^{13}\right)\]}
\begin{equation}\tag*{[20]}M\surd(1-kk)=\frac{k\text{ Numerator pro }\frac1k}{C}\quad\text{demonstr. GALEN}\end{equation}
\par\vfill\hfill24*
\clearpage
\SourceLeaf{195}\ArithmeticHead{188}{L}{ANALYSIS.\quad NACHLASS.}
\begin{equation}\tag*{[21]}\text{Obser[vatio:]}\qquad\frac{M\sin75^0}{M\sin15^0}=\surd3\end{equation}
{\small\begin{equation}\tag*{[22]}\left\{\begin{array}{rrrrl}
1&-4&+32&-176&\\
&&-12&+192&1-4zz+20z^4-64z^6\ldots=\text{Periph. Ellips.}\\
&&&-80&
\end{array}\right.\end{equation}}
\NumberedItem{5}
\noindent [S. 12]
\begin{equation}\tag*{[24]}\left\{\begin{gathered}
\text{Posito termino constante expressionis}\\
(1+\mu\cos\varphi^2)^{-\frac12}\\
=A,\text{ erit}\\
M\surd(1+\mu)=\frac1A
\end{gathered}\right.\end{equation}
\begin{equation}\tag*{[24]}\left\{\begin{gathered}
\text{Terminus constans expressionis }\left(hh-2\cos\varphi+\frac1{hh}\right)^{-\frac12}\\
\text{est}=\frac1{\text{Terminus medius inter }h-\frac1h\text{ et }h+\frac1h}
\end{gathered}\right.\end{equation}
\begin{equation}\tag*{[25]}\begin{gathered}
\text{Si }(\mu\mu-\cos\varphi^2)^{-\frac12}\text{ in seriem}\\
A+2B\cos2\varphi+\text{etc.}\\
\text{evolvi concipitur atque ponitur}\\
A=F\mu,\ B=G\mu
\end{gathered}\end{equation}
\noindent erit
\begin{equation}\tag*{[26]}\left\{\begin{aligned}
F\frac12\left(\mu+\frac1\mu\right)&=2\mu F\mu\mu\\
G\frac12\left(\mu+\frac1\mu\right)&=\frac1\mu G\mu\mu+\frac1\mu F\mu\mu
\end{aligned}\right.\end{equation}
\par\indent Ponendo
\begin{equation}\tag*{[27]}\frac{Gx}{Fx}=Hx\qquad x,\ x',\ x''\text{ etc.}\end{equation}
\noindent erit
\begin{equation}\tag*{[28]}\begin{aligned}
Hx&=\frac1{2x'}+\frac1{4x'x''}+\frac1{8x'x''x'''}\text{ etc.}\\
&=\frac18x^{-2}+\frac1{16}x^{-4}+\frac{41}{8024}x^{-6}\text{ etc.}
\end{aligned}\end{equation}
\[Hxx=2xxH\frac12\left(x+\frac1x\right)-1\]
\clearpage
\endgroup

